\documentclass[10pt,a4paper]{amsart}
\usepackage[margin=19mm]{geometry}
\usepackage{amsmath,amssymb,mathtools}
\usepackage[scr=rsfs,cal=euler]{mathalfa}
\usepackage{xfrac}
\usepackage[unicode,hidelinks,bookmarks=false]{hyperref}
\newcommand{\ie}{i.e.\ }
\newcommand{\cf}{cf.\ }
\newcommand{\resp}{respectively\ }
\newcommand{\Subsection}{\subsection}
\providecommand{\nobreakdash}{\nobreak\hbox{-}}
\setlength{\emergencystretch}{2em}
\multlinegap0pt
\usepackage{bm}
\usepackage{bbm}
\usepackage{dsfont}
\usepackage{xspace}
\usepackage{cancel}
\usepackage{mathtools}
\usepackage{amsthm}
\makeatletter
\@ifundefined{thm}{\newtheorem{thm}{Theorem}}{}
\makeatother
\makeatletter
\newcommand{\PV}{\texttt{PV}}
\newcommand{\R}{\mathbb R}
\@ifundefined{proof}{\newenvironment{proof}[1][Proof]{\textbf{#1.} }{\ \rule{0.5em}{0.5em}}}{}
\makeatother
\title[Existence, asymptotic behaviour and convergence of a general]{Existence, asymptotic behaviour and convergence of a generalised 3D Muskat problem in stable regime}
\author{Qasim Khan, Anthony Suen, Bao Quoc Tang}
\date{Source: arXiv:2603.16172v2 (CC BY 4.0)}
\begin{document}
\maketitle

\noindent\emph{Source and licence.} Existence, asymptotic behaviour and convergence of a generalised 3D Muskat problem in stable regime. Version arXiv:2603.16172v2 (CC BY 4.0), distributed under the Creative Commons Attribution 4.0 International licence, as recorded in the arXiv licence metadata for this exact version and shown on the abstract page https://arxiv.org/abs/2603.16172v2. Full rights notice: licenses/arxiv-2603.16172-CC-BY-4.0.txt.\par\smallskip

\noindent\textit{Editorial scope.} Counted unit: the Thm whose source label is \texttt{L2 convergence thm} in the article (source0.tex lines 1170--1175, source label \texttt{L2 convergence thm}), delivered together with the article's own complete original proof of it (source0.tex lines 1176--1255, one \texttt{proof} environment of 80 lines). No mathematics is written, repaired or re-derived: statement and proof are the article's own text line for line. Required context retained uncounted: contour eqn (lines 259--264); local existence thm (lines 300--302); def of M alpha (f) (lines 377--379); uniform bound on H4 norm (lines 501--503). Every definition, display and label of those lines is retained unchanged. Omitted from this package and not redistributed: 1--258, 265--299, 303--376, 380--500, 504--1169, 1256--1406. Nothing inside a retained excerpt is omitted or altered: every retained body is delivered exactly as the article's lines are written, so no omission receipt of any kind accompanies this package. Citations: the retained excerpts carry no citation command at all, so no bibliography entry of the article is transcribed and no reference is redistributed. Reference adaptation: none was needed and none was made; every reference inside the delivered unit resolves inside it, so no unresolved reference or citation is produced. Printed slips kept literal: no printed word, symbol, hypothesis or number of the counted statement or of its proof was altered. Layout and class: the article's own class is \texttt{amsart} with packages including verbatim, xcolor, graphicx, amssymb, amsmath, amsthm, amsfonts, float, hyperref, mathrsfs, enumitem, scalerel, stackengine; the wrapper is the lane's pinned \texttt{amsart} editorial wrapper at 10pt on A4, which restates the theorem-like environments the retained text needs and adds the article's own macro definitions that the retained text needs (\texttt{\textbackslash PV}, \texttt{\textbackslash R}), with extra wrapper packages (bm, bbm, dsfont, xspace, cancel, mathtools). The delivered numbering is the unit's own. Assets: no retained span contains a figure environment, a graphics-inclusion command, a PDF-page inclusion, a table or a TikZ picture; no image file is copied into this package, the package declares no asset dependency and no image file is redistributed with it. Licence: version arXiv:2603.16172v2 of this article is distributed under the Creative Commons Attribution 4.0 International licence, as recorded in the arXiv licence metadata for this exact version and shown on the abstract page https://arxiv.org/abs/2603.16172v2. A full rights notice is delivered at licenses/arxiv-2603.16172-CC-BY-4.0.txt. Characters: every retained excerpt is pure ASCII, so the delivered engine, XeLaTeX with its default Latin Modern text font, reports no missing character.
\begin{equation}\label{contour eqn}
\begin{split}
\frac{\partial f(x,t)}{\partial t} &= \frac{\rho_2 - \rho_1}{4\pi} \, \PV \, \int_{\mathbb{R}^2} \frac{ \left( \nabla f(x,t) - \nabla f(x - y, t) \right) \cdot y }{ \left[ |y|^2 + (f(x,t) - f(x - y, t))^2 \right]^{\frac{3 + \alpha}{2}}} \, dy, \\
f(x,0) &= f_0(x), \quad x=(x_1, x_2) \in \mathbb{R}^2.
\end{split}
\end{equation}

\begin{thm}\label{local existence thm}
Let $f_0\in H^k(\R^2)$ for $k\ge4$. Then for any $\alpha\in[0,1)$, there exists a time $T>0$ such that the contour equation \eqref{contour eqn} possess a unique solution in $C^1([0,T];H^k(\R^2))$ with $f(x,0)=f_0(x)$.
\end{thm}

\begin{equation}\label{def of M alpha (f)}
M_\alpha(f) = \max_x\Big| \PV \int_{|y|<1}  \frac{\left(f(x)-f(x-y)\right)\cdot \left(\nabla f(x)-\nabla f(x-y)\right)\cdot y}{\left[ |y|^2 + \left(f(x) - f(x - y)\right)^2 \right]^{\frac{5 + \alpha}{2}}}dy\Big|.
\end{equation}

\begin{align}\label{uniform bound on H4 norm}
\sup_{0\le \alpha<\bar{\alpha}}\sup_{0\le t\le \bar{T}}\|f_\alpha\|_{H^4}(t)\le \bar{M}.
\end{align}

\begin{thm}\label{L2 convergence thm}
For $\alpha\in[0,\frac{1}{2})$, $k\ge4$ and $f_0\in H^k(\R^2)$, assume that $f_\alpha$ and $f$ are solutions in $C^1([0,\bar{T}];H^k(\R^2))$ as defined above. Then for all $t\in[0,\bar{T}]$, we have
\begin{align}\label{L2 convergence}
\lim_{\alpha\to0^+}\|f_\alpha - f\|_{L^2}(t)=0.
\end{align}
\end{thm}

\begin{proof}
Throughout this proof, $C_{\bar{M}}$ always denotes a positive constant that depends on $\bar{M}$ given in \eqref{uniform bound on H4 norm} but is independent of $\alpha$. 

Let $g_{\alpha}=f_{\alpha}(x, t)-f(x, t)$, then $g_{\alpha}$ satisfies
\begin{equation}\label{egn for g alpha}
\begin{split}
\frac{\partial g_{\alpha}}{\partial t}(x, t) & = \PV \int_{\mathbb{R}^{2}} \frac{\left(\nabla g_{\alpha}(x, t)-\nabla g_{\alpha}(x-y, t)\right) \cdot y}{A_{\alpha}^{\frac{3+\alpha}{2}}} d y \\
& \qquad+ \PV \int_{\mathbb{R}^{2}} \frac{(\nabla f(x, t)-\nabla f(x-y, t)) \cdot y}{A_{\alpha}^{\frac{3+\alpha}{2}}} d y \\
& \qquad- \PV \int_{\mathbb{R}^{2}} \frac{(\nabla f(x, t)-\nabla f(x-y, t)) \cdot y}{A_{0}^{\frac{3}{2}}} d y. 
\end{split}
\end{equation} 
Multiplying \eqref{egn for g alpha} by $g_{\alpha}$ and integrating with respect to $x$,
\begin{equation*} 
\begin{split}
 \frac{\partial}{\partial t} \int_{\mathbb{R}^{2}}\left|g_{\alpha}(x, t)\right|^{2} d x&=\PV \int_{\mathbb{R}^{2}} 2g_{\alpha}(x, t) \int_{\mathbb{R}^{2}} \frac{\left(\nabla g_{\alpha}(x, t)-\nabla g_{\alpha}(x-y, t)\right) \cdot y}{A_{\alpha}^{\frac{3+\alpha}{2}}} d y d x \\
&\,\,\,+\PV \int_{\mathbb{R}^{2}} 2g_{\alpha} \int_{\mathbb{R}^{2}}(\nabla f(x, t)-\nabla f(x-y, t)) \cdot y \cdot\left(\frac{1}{A_{\alpha}^{\frac{3+\alpha}{2}}}-\frac{1}{A_{\alpha}^\frac{3}{2}}\right) \\
&\,\,\,+\PV \int_{\mathbb{R}^{2}} 2g_{\alpha} \int_{\mathbb{R}^{2}}(\nabla f(x, t)-\nabla f(x-y, t)) \cdot y \cdot\left(\frac{1}{A_{\alpha}^\frac{3}{2}}-\frac{1}{A_{0}^\frac{3}{2}}\right) \\
&\triangleq \mathbf{E}_{1}+\mathbf{E}_{2}+\mathbf{E}_{3}. 
\end{split}
\end{equation*}
We estimate the terms $\mathbf{E}_{1}$, $\mathbf{E}_{2}$ and $\mathbf{E}_{3}$ as follows. First of all, we rewrite $\mathbf{E}_{1}$ as follows (we drop the variable $t$ for simplicity):
\begin{align*}
\mathbf{E}_{1}=\mathbf{E}_{1,1}+\mathbf{E}_{1,2}+\mathbf{E}_{1,3},
\end{align*}
where
\begin{align*}
\mathbf{E}_{1,1}&=\PV\int_{\mathbb{R}^{2}} 2g_{\alpha}(x) \int_{\mathbb{R}^{2}} \frac{\nabla g_{\alpha}(x) \cdot y}{A_{\alpha}^{\frac{3+\alpha}{ 2}} }\\
\mathbf{E}_{1,2}&=-\PV\int_{\mathbb{R}^{2}} 2g_{\alpha}(x) \int_{\mathbb{R}^{2}} \frac{g_{\alpha}(x)-g_{\alpha}(y)}{A_{\alpha}^{\frac{3+\alpha}{2}}} \\
\mathbf{E}_{1,3}&=\PV\int_{\mathbb{R}^{2}} 2(3+\alpha)g_{\alpha}(y) \int_{\mathbb{R}^{2}} (g_{\alpha}(x)-g_{\alpha}(y))\\
&\qquad\qquad\qquad\qquad\times\left[\frac{\left(f_{\alpha}(x)-f_{\alpha}(y)\right) \cdot\left(f_{\alpha}(x)-f_{\alpha}(y)-\nabla f_{\alpha}(y) \cdot(x-y)\right)}{A_{\alpha}^{\frac{5+\alpha}{2}}}\right]. 
\end{align*}
Upon integrating by parts, we have
\begin{align*}
\mathbf{E}_{1,1}&=\PV \int_{\mathbb{R}^{2}} (3+\alpha)|g_{\alpha}(x, t)|^{2}\int_{\mathbb{R}^{2}}\frac{\left(f_{\alpha}(x)-f_{\alpha}(x-y)\right)\left(\nabla f_{\alpha}(x)-\nabla f_{\alpha}(x-y)\right) \cdot y}{A_{\alpha}^{\frac{5+\alpha}{ 2}} }\\
&\le (3+\alpha)\left\|f_{\alpha}\right\|_{ C^{1}}\left\|g_{\alpha}\right\|_{L^{2}}^{2}+(3+\alpha) M_{\alpha}\left(f_{\alpha}\right)\left\|g_{\alpha}\right\|_{L^{2}}^{2}\\
&\le 4\left\|f_{\alpha}\right\|_{ C^{1}}\left\|g_{\alpha}\right\|_{L^{2}}^{2} + 4M_{\alpha}\left(f_{\alpha}\right)\left\|g_{\alpha}\right\|_{L^{2}}^{2}
\end{align*}
where $M_\alpha(\cdot)$ is given by \eqref{def of M alpha (f)}. Following the proof of Theorem~\ref{local existence thm}, Using the uniform bound \eqref{uniform bound on H4 norm}, we readily have
\begin{align*}
4\left\|f_{\alpha}\right\|_{ C^{1}}+4M_{\alpha}(f_{\alpha})\le C_{\bar{M}},
\end{align*}
and hence
\begin{align*}
\mathbf{E}_{1,1}\le C_{\bar{M}}\left\|g_{\alpha}\right\|_{L^{2}}^{2}.
\end{align*}
On the other hand, by making change of variables,
\begin{align*}
\mathbf{E}_{1,2}=-\frac{1}{2}\int_{\R^2}\int_{\R^2}\frac{(g_\alpha(x)-g_\alpha(y))^2}{A_{0}^\frac{3+\alpha}{2}}\le 0,
\end{align*}
while for $\mathbf{E}_{1,3}$, following the similar method for estimating $\mathbf{J}_3$ as in the proof of Theorem~\ref{local existence thm}, we have
\begin{align*}
\mathbf{E}_{1,3}\le 8\left[\left\|f_{\alpha}\right\|_{C^{2, \delta}}^{4}+\left\|f_{\alpha}\right\|_{C^{1}}\|f\|_{C^{2, \delta}}+1\right]\left\|g_{\alpha}\right\|_{L^{2}}^{2}\le C_{\bar{M}}\left\|g_{\alpha}\right\|_{L^{2}}^{2}.
\end{align*} 
Therefore we obtain
\begin{align}\label{bound on E1}
\mathbf{E}_{1}\le C_{\bar{M}}\left\|g_{\alpha}\right\|_{L^{2}}^{2}.
\end{align}
Next for the term $\mathbf{E}_2$, using mean value theorem and the uniform bound \eqref{uniform bound on H4 norm}, we have
\begin{align*}
\left|\frac{1}{A_{\alpha}^{\frac{3+\alpha}{2}}}-\frac{1}{A_{\alpha}^{3 / 2}}\right|&\le\left|-\frac{3+\alpha}{2}+\frac{3}{2}\right|\sup_{r\in(-\frac{3+\alpha}{2},-\frac{3}{2})}|A_{\alpha}^r||\log(A_\alpha)|\\
&\le\frac{\alpha}{2}|y|^{-3}\log(M^2)=\alpha|y|^{-3}\log(M).
\end{align*}
Therefore by applying the above estimate to $\mathbf{E}_2$, it implies
\begin{align}\label{bound on E2}
\mathbf{E}_2\le \alpha C_{\bar{M}}\left\|g_{\alpha}\right\|_{L^{2}}.
\end{align}
For the term $\mathbf{E}_3$, using mean value theorem and the uniform bound \eqref{uniform bound on H4 norm}, we also have
\begin{align*}
|A_{0}^\frac{3}{2}-A_{\alpha}^\frac{3}{2}|\le C_{\bar{M}}|g_{\alpha}(x)-g_{\alpha}(x-y)||y|^2,
\end{align*}
and hence
\begin{align}\label{bound on E3}
\mathbf{E}_3\le C_{\bar{M}}\left\|g_{\alpha}\right\|_{L^{2}}^{2}.
\end{align}
Combining \eqref{bound on E1}, \eqref{bound on E2} and \eqref{bound on E3}, we obtain
\begin{align}\label{full bound on g alpha}
\frac{\partial}{\partial t}\|g_{\alpha}\|^2_{L^2}\le 2C_{\bar{M}}\|g_{\alpha}\|^2_{L^2}+\alpha C_{\bar{M}}\|g_{\alpha}\|_{L^2}.
\end{align}
Applying Gr\"{o}nwall's inequality to \eqref{full bound on g alpha} and recalling that $g_{\alpha}(x,0)=0$, the results \eqref{L2 convergence} follows immediately by taking $\alpha\to0^{+}$.
\end{proof}

\end{document}
